\documentclass[12pt]{article}
\usepackage[utf8]{inputenc}
\usepackage[T1]{fontenc}
\usepackage{hyperref}
\title{Subi Language Resources for K-12 Curriculum Development}
\author{Compiled by Claude Sonnet 5.5 (claude-sonnet-5-5) for Andr\'as Kornai}
\date{}
\begin{document}
\maketitle

\section{Introduction}
Subi (Kisubi; ISO 639-3 xsj) is a Bantu language of Tanzania. It is spoken by the Abhasubi people, mainly in Biharamulo District, Kagera Region. Sources give about 105,000 to 106,000 speakers. It should not be confused with Shubi. The Living Dictionary describes it as not yet endangered, but under pressure from Kiswahili. Very little published educational material was found in search results.

\section{Primary Educational Resources}
No primers, textbooks or school curricula in Subi were found in search results. Nothing is known of Subi being taught in schools.

\section{Linguistic and Reference Materials}
\begin{itemize}
\item Subi Living Dictionary, a community-oriented dictionary by Cyprian Vumilia and John Gluckman: \url{https://livingdictionaries.app/subi} (about page: \url{https://livingdictionaries.app/subi/about})
\item Wikipedia, Subi language: \url{https://en.wikipedia.org/wiki/Subi_language}
\item Glottolog entry: \url{https://glottolog.org/resource/languoid/iso/suj}
\end{itemize}

\section{Religious and Community Resources}
\begin{itemize}
\item Joshua Project profile of the Subi people (describes them as mainly Muslim, centred on Biharamulo District): \url{https://joshuaproject.net/people_groups/21582}
\end{itemize}
No Subi scripture or community organisation was found.

\section{Digital and Audio Resources}
\begin{itemize}
\item Global Recordings Network language page, which lists the language but reports no recordings available: \url{https://globalrecordings.net/en/language/xsj}
\end{itemize}

\section{Curriculum Development Context}
The documented resources are a single dictionary and general reference pages. A Subi curriculum would need to be built largely from scratch. The first steps would be an agreed orthography and community-led materials, using the Living Dictionary as a lexical base. Kiswahili is the school language in Tanzania, and instruction in Subi would have to be arranged with the community.

\end{document}